\begin{abstract}
不完整的来源观测会留下多个读自关系与修改序补全，分析任务因而需要回答条件下哪些来源和写序仍然可行，而不仅是寻找一个见证。困难在于：每个完成的执行图具有小接口，并不保证所有候选的共同表示也具有小接口。本文研究固定事件的 Release--Acquire 观测，提出由协议共同因果关系定义的未决前沿。对于前沿以外已经确定先于未来相关访问的历史，每个地址只需保留少量极大写；结合线程尾事件，可构造覆盖全部 RF／CO 候选的连续处理计划。我们从单生产者／单消费者缓冲的容量反馈及重置／唤醒屏障推导统一证书，并证明精确编译对全部原始赋值保持投影，使已退休事件的来源身份仍可用于条件支持与假设撤销。实验中，容量为三的分支缓冲从四轮增至三十二轮时，最大可达层状态保持为三十；容量增至六则触及当前状态预算。组件分解和日志快照频率又会改变编译与 SAT 的收益关系。这些结果说明，协议结构能够支持可复用的精确分析，同时给出其计算收益的具体适用条件。
\end{abstract}

\section{引言}
一条并发日志记录“读到了值 $1$”，通常没有记录它来自哪次写入。如果同一线程多次发布值 $1$，早、晚两次发布可能对应不同的同步知识：早发布发生在某次数据写之前，晚发布发生在它之后。加入“此次读来自晚发布”的假设，可能迫使另一个地址上的写序方向；撤销这条来源记录，先前排除的方向又应恢复。定位不一致日志、枚举一个字段的替代解释，以及比较候选来源，都需要反复回答这类条件问题。

一次可满足性求解提供一个完整见证，无法直接说明其余来源是否也可行。对每个新条件重新搜索全部支持虽然精确，却可能重复处理同一组同步约束。将固定观测的完成空间预编译为可复用表示，是一种自然选择。它的主要障碍不是如何在图上做一次可达性，而是如何在尚未选择读源和修改序时，控制要表示的整个候选空间。

这项困难具有明确的量词差异。为每个完成分别寻找一个小分解，允许布局依赖已选来源；查询表示却需要在未来假设未知时保留所有合法选择。同一批同值写和读可以使每个完成的 RF 图很窄，而所有候选取并后形成完全二部图。已有基于完成图的结构保证因而不能直接作为统一编译接口的保证。

本文利用同步协议所建立的\emph{共同因果关系}。即使一个读的具体来源未定，所有候选来源仍可能具有同一个前序。随着编译推进，某些历史事件已确定先于未来每个相关访问。对于这部分稳定历史，一个旧写若被同地址的更晚稳定写覆盖，就不可能再作为未来读的来源，也不可能是未来写的固定修改序覆盖前驱。剩下的稳定极大写，每个地址至多有线程数那么多个。必须继续显式保留的其他事件集中在未决因果前沿中。

由此得到本文的结构定理：若未决前沿至多有 $c$ 个事件，线程数为 $p$，地址数为 $a$，则整个 RF／CO 候选图有可直接构造的处理计划，最大事件袋大小至多 $c+ap+p+a+1$。这里界定的是所有候选共同使用的接口，允许一条 RF 或 CO 覆盖边跨过任意长历史。事件退出接口后，其关系作用由联合摘要保存；它曾选择的原始文字仍保留在编译路径上。安全遗忘关系状态与保留历史查询身份因而可以同时实现。

结构参数还必须能从程序纪律推出。我们给出两种具体来源。双模索引缓冲的容量反馈限制所有尝试的来源窗口，上界通过真实前缀中的因果环反证得到，无需等待完整日志的末端。星形汇聚／唤醒屏障则利用重置写推动到达来源前进，再经反馈锁定当前调用；适当布局使任意多失败轮询不增加未决前沿。两个协议构造保留完整有限前缀，并为尚未选定的 RF／CO 提供同一证书和计划。

本文的贡献集中于三个联系。第一，证明共同因果前沿与稳定极大写如何约束整个候选图，并明确二元选择及连续生存期所需条件。第二，给出反馈缓冲和屏障的可检查协议推导，以及失败读链和地址封闭组合的构造，从程序参数获得接口界。第三，将这一结构接到精确原始投影和完整支持查询上，以实现及冻结实验说明反复查询的收益何时出现、何时被组件缓存或构造成本抵消。通用关系分析、自动机检查、知识编译及单字段修复分别作为这些联系的基础和使用方式。

实验中的正面结果与边界同样影响结论。固定容量的缓冲在所测历史增长下保持小的可达层状态，并保留跨阶段来源关联；更大容量会使联合状态迅速增加。屏障中的可提取组件让局部缓存有效应对大量新全局条件，而实际日志诊断中，快照频率足以改变编译与 SAT 的端到端排名。本文据此研究一种结构可解释的支持编译方法，而不以普遍性能领先作为理论前提。
